\section{Conclusion}
\label{sec:conclusion}

\ours shows that three simple ideas, importance-scored keyframe memory, temporal query routing, and a two-stage perception-synthesis pipeline, are enough to bring video Q\&A into the streaming regime. None of the individual components is exotic. Their value comes from how they fit together: the SKM captures what matters, the TQR tells the generator where to look, and the BLIP + Flan-T5 pipeline produces coherent answers with low latency on a GPU.

On LiveQA-Bench, \ours achieves 39.7\% overall with 59.3\% on instant queries, showing that temporal routing and importance-scored memory work together effectively.
The 4.8-point gain from $N\!=\!16$ to $N\!=\!64$ confirms that larger memory benefits historical queries in particular.
All of this works with frozen pre-trained checkpoints and no end-to-end training.
Full results on NExT-QA, EgoSchema, OVO-Bench, and Ego4D-NLQ under the streaming protocol, along with baseline comparisons, will be completed for the final version.

\paragraph{Limitations.}
The SKM scores frames by visual similarity in CLIP embedding space. Events that matter semantically but look ordinary on camera, a quiet conversation, a subtle hand gesture, receive low scores and risk eviction.
The TQR bins questions into three discrete scopes via keyword matching. Requests like ``between minutes 3 and 7'' fall outside its vocabulary. A learned scope classifier would cover more ground.
The scope boundaries are fixed: the recency window $W\!=\!30$\,s does not adapt to the event density of the stream, which may cause scope misalignment in very slow or very fast-paced videos (\cref{sec:method}).
The system handles only one camera stream. We do not address multi-view fusion.

\paragraph{Ethics.}
A system that watches live video and answers arbitrary questions invites misuse for surveillance. Operators should deploy \ours only with informed consent from everyone on camera. The demo processes video locally and stores no frames beyond the active session.

\paragraph{Future work.}
Four directions seem most promising. First, replacing keyword-based routing with a small trained classifier would handle ambiguous questions better and adapt to varying event densities. Second, accepting audio alongside video frames would give the system a richer picture of the scene. Third, moving from one camera to many, as in a smart-home or retail setting, would test whether the memory and routing design scales. Fourth, integrating the temporal routing mechanism into recent streaming LLMs such as VideoLLM-online~\cite{chen2024videollmonline} or Dispider~\cite{he2024dispider} could combine the strengths of both paradigms.
